% Capítulo IV — Construcción e Implementación del Sistema (por incrementos)
% No redactar un incremento hasta que su código exista (documento/04-pendientes-preparacion.md).
\capitulo{4}{IV}{Construcción e Implementación del Sistema}

\subsection{INCREMENTO 1: PROCESOS DEL PORTAL ANTERIOR MIGRADOS AL ERP}

\subsection{INCREMENTO 2: AUTENTICACIÓN CENTRALIZADA CON KEYCLOAK}

\subsection{INCREMENTO 3: PIPELINE DE DATOS Y ENTRENAMIENTO DE MODELOS}
\subsubsection{INTEGRACIÓN DUCKDB CON MYSQL Y POSTGRESQL}
\subsubsection{TRANSFORMACIÓN DE DATOS CON POLARS Y FEATURE ENGINEERING}
\subsubsection{MODELO DE DESVIACIÓN EN CERTIFICACIONES}
\subsubsection{MODELO DE SOBRECOSTO DE PROYECTOS}
\subsubsection{MODELO DE RETRASO EN LA FECHA DE CIERRE}

\subsection{INCREMENTO 4: MICROSERVICIO FASTAPI E INTEGRACIÓN CON EL ERP}
\subsubsection{DESARROLLO DEL MICROSERVICIO FASTAPI}
\subsubsection{ENDPOINTS DE PREDICCIÓN EN NESTJS}
\subsubsection{DASHBOARD DE INTELIGENCIA PREDICTIVA EN ANGULAR}

\subsection{INCREMENTO 5: INTEGRACIÓN CON POWER BI Y MCP}
\subsubsection{API DE DATOS PARA POWER BI}
\subsubsection{SERVIDOR MCP}

\subsection{EVALUACIÓN Y VALIDACIÓN}
\subsubsection{MÉTRICAS DE RENDIMIENTO DE MODELOS}
\subsubsection{PRUEBAS CON USUARIOS CLAVE}
